\subsection{A general left-identity extension of the five-element band}

\begin{theorem}[Proved: constant actors with a common left identity]
\label{CP2:A-constant-actor-left-identity-quotient}
Let $M$ be a finite monoid with minimal ideal the
$2\times2$ rectangular band $K=\{e_{ij}\}$.
Put $P=\{1\}\sqcup K$ and $V=M\setminus P$.
Suppose $V$ is nonempty and there exist $e\in K$
and $q\in V$ such that
\[
 vk=ek,\quad kv=ke\quad(v\in V,\ k\in K),
 \qquad qv=v\quad(v\in V).                                    \tag{1}
\]
Then the retraction $\rho:M\to P$ that sends every
$v\in V$ to $e$ has the following property.
For every finite nonempty left-$N$, right-$M$ biset
$S$, it extends to an actual quotient of collages
$\mathcal C_S\to\mathcal C_U$ inducing rational
homology and cohomology isomorphisms through degree two.
For $N=R_3$, all rational degree-two products of
$B\mathcal C_S$ vanish, without any state-count,
Betti-number, or boundedness assumption.
\end{theorem}
\begin{proof}
Relabel the band so that $e=e_{00}$.
Equation (1) implies $ve=ev=e$ for every $v\in V$.
A product of two elements of $V$ cannot be the global
identity: otherwise they would be units in the finite
monoid, whereas their constant left action on the two
band rows is not a permutation.
If their product lies in $K$, its left and right band
actions are those of $e$, so the product is exactly $e$.
Thus $V\cup\{e\}$ is closed, and sending it to $e$
does define the stated monoid retraction.
The element $q$ is idempotent and $qe=eq=e$.

Write $R=\mathbb QM$, $B=\mathbb QP$,
$h=q-e$, and $p=1-h$.  Then $h,p$ are orthogonal
idempotents.  The band ideal is annihilated on both
sides by $h$.
For $v\in V$, the left-identity condition gives
\[
 pv=e,\qquad vp=v-vq+e,\qquad vh=vq-e.
\]
Since $vq$ belongs to $V\cup\{e\}$, it is fixed
on the left by $q$.  Therefore $pRh=0$ and
\[
 H=hRh=Rh,\qquad pRp\cong B,
 \qquad W=hRp=\mathbb Q\{v-vq:v\in V\}.                       \tag{2}
\]
The isomorphism $pRp\to B$ is the restriction of
$\rho$; its basis is $p$ together with the four
band elements.
Let $\Delta$ be the right-$B$ character on which
the identity acts as one and every band actor as zero.
The right action on $W$ is exactly a direct sum of
copies of $\Delta$, because $WK=0$ and $Wp=W$.
There is an exact sequence of left-$R$, right-$B$
bimodules
\[
 0\longrightarrow W\longrightarrow Rp
      \longrightarrow B\longrightarrow0.                     \tag{3}
\]

Choose any nonnegative free left-$H$ resolution
$Q_*\to W$, retaining the external right-$B$ action
$\Delta$ on all its terms.
The left-$R$ module $Rh$ is projective and, by (2),
is the regular free right-$H$ module.
Thus $Rh\otimes_H Q_*$ resolves $W$ by left-$R$
projectives.  Splice it into (3).
For $L=\mathbb Q[S]$, the resulting tensor complex
$F$ computes $L\otimes_R^{\mathbf L}B$.
It has $F_0=Lp$ and, in every positive degree,
right-$B$ action $\Delta$.
All these constructions retain the commuting
left-$\mathbb QN$ action.

Let $C,E$ be the actors of $q,e$ on $S$.
First identify $C(s)$ with $E(s)$ for every state.
The resulting set $T$ is left-$N$, right-$P$ equivariant,
and the explicit integral basis
\[
 s-C(s)+E(s)\quad(s\notin CS),\qquad t\quad(t\in ES)
\]
identifies $Lp$ with $\mathbb Q[T]$.
The module $Lh$ has basis $t-E(t)$ with
$t\in CS\setminus ES$.
For every $v\in V$, the product of such a vector
with the generator $v-vq$ of $W$ maps under this
identification to
\[
 [tv]-[E(t)].                                                   \tag{4}
\]
Indeed its literal difference is $tv-tvq$, and
$tvq=C(tv)$ was identified with
$E(tv)=E(t)$ in the first quotient.

Let $U$ be the component set of the graph on $T$
with all edges (4).
The image of $F_1\to F_0$ is exactly the incidence
span of these edges, since the chosen free resolution
surjects onto $W$ and tensor products are right exact.
Consequently $H_0(F)=\mathbb Q[U]$ with its literal
state basis.
The relations are left-$N$ stable by commutation,
and right-$P$ stable because each edge's endpoints
have equal images under every band actor.
Moreover $qv=v$ says that every $v$-image of $S$
already has a representative $tv$ with $t\in CS$.
Hence every added actor descends to $e$ on $U$.
This proves that the quotient is the actual
$(N,\rho)$-equivariant map $S\to U$.

The natural comparison $F\to\mathbb Q[U]$ is the
derived adjoint of this actual state map: in degree
zero it sends $lp$ to the image of $l$ because
$\rho(p)=1$.
Its kernel complex is degreewise a right-$B$ module
of type $\Delta$ and has zero degree-zero homology.
For degree zero this uses the incidence image
$LhW\subseteq Lp$, which is annihilated by $K$;
for positive degrees it follows from the construction
of $Q_*$.
The natural good truncation therefore identifies
the kernel up to quasi-isomorphism with
$G\otimes_{\mathbb Q}\Delta[1]$ for a nonnegative
left-$\mathbb QN$ complex $G$.
No splitting of left-$N$ modules is required.

As in Theorem~\ref{CP2:A-eight-triangular-low-degree-quotients},
the band resolution gives
$\Delta\otimes_B^{\mathbf L}\mathbb Q\simeq\mathbb Q[1]$.
The actual bridge-comparison fiber is consequently
$(\mathbb Q\otimes_{\mathbb QN}^{\mathbf L}G)[2]$,
which has no homology below degree two.
The trivial right-$R$ module annihilates $Rh$, so
the same resolution shows that the actual corner map
$BM\to BP$ is a rational homology equivalence.
The two-block mapping-cone sequence shifts the bridge
fiber twice more.  The collage comparison cone has
no homology below degree four, proving the asserted
isomorphisms through degree two.

For $N=R_3$, the target collage has homology bounded
in degree four by the length-two $P$ resolution and
length-one right-$N$ resolution.
The known all-state right-zero-three graph theorem
therefore gives zero degree-two products on that
target.  Every degree-two class on the original
collage is its pullback, so its products vanish as well.
\end{proof}

\begin{remark*}
The two explicit order-eight tables treated immediately
before this theorem correspond to $q=c$.
The general argument allows an arbitrary finite algebra
$H$ and an arbitrary left-$H$ module $W$.
It uses only the lower degree of the comparison fiber;
it neither assumes nor discards its possible unbounded
homology.
\end{remark*}
